\section{Tensor Schwinger functions}\label{sec:tensors}
The following estimate is the argument of
\cite[Section~5.3]{QED3Manuscript}. It applies to any
Gaussian-dominated real two-form law in three Euclidean dimensions.
For \(n\ge1\), set
\[
 \mathcal S_n=\mathcal S((\R^3)^n;(\Lambda^2\C^3)^{\otimes n}),\qquad
 q_m(\Psi)=\max_{\substack{|\alpha|\le m\\\text{components}}}
 \sup_y(1+|y|^2)^{m/2}|\partial^\alpha\Psi(y)|.
\]
\begin{proposition}\label{prop:tensor-growth}
Let \(P\) be a law of real tempered two-forms with
\(\E_Pe^{F(f)}\le e^{\|f\|_2^2/2}\) for all real Schwartz \(f\).
For each \(n\ge1\) there is a continuous linear map
\(\Psi\mapsto Q_\Psi\) from \(\mathcal S_n\) to \(L^2(P)\)
which agrees with \(\prod_aF(f_a)\) on decomposable tests.
The functional \(S_n(\Psi)=\E_PQ_\Psi\) is the unique continuous
extension of the smeared moments to \(\mathcal S_n\), and
\begin{equation}\label{eq:tensor-growth}
 |S_n(\Psi)|\le2\,66^n\sqrt{n!}\,q_{6n}(\Psi).
\end{equation}
\end{proposition}

\begin{proof}
Let \(h_\kappa=h_ke_{\ell m}\) be the real Hermite two-form basis.
The eigenvalues of \(A=1+|x|^2-\Delta\) are
\(\lambda_\kappa=2|k|_1+4\). Expand in Schwartz space
\[
 \Psi=\sum_{\boldsymbol\kappa}c_{\boldsymbol\kappa}(\Psi)
 h_{\kappa_1}\otimes\cdots\otimes h_{\kappa_n}.
\]
Counting eigenvalues gives
\[
 Z_4:=\sum_\kappa\lambda_\kappa^{-4}
 =\frac3{32}\sum_{m\ge0}\frac{m+1}{(m+2)^3}
 \le\frac3{32}\left(\frac{\pi^2}{6}-1\right)<\frac1{16}.
\]
Weighted Cauchy--Schwarz and Parseval imply
\[
 \sum_{\boldsymbol\kappa}|c_{\boldsymbol\kappa}(\Psi)|
 \le Z_4^{n/2}\|(A^2)^{\otimes n}\Psi\|_2.
\]
For \(M(x)=1+|x|^2\),
\[
 A^2=M^2-2M\Delta-4x\cdot\nabla-6+\Delta^2.
\]
The sum of the absolute differential coefficients is bounded by
\(34M^2\). Applying this in each argument gives
\[
 |(A^2)^{\otimes n}\Psi(x)|
 \le34^nq_{6n}(\Psi)\prod_{a=1}^n(1+|x_a|^2)^{-1}
\]
componentwise. Since
\(\int_{\R^3}(1+|x|^2)^{-2}\,\dd x=\pi^2\) and there are
\(3^n\) components,
\begin{equation}\label{eq:Hermite-l1}
 \sum_{\boldsymbol\kappa}|c_{\boldsymbol\kappa}(\Psi)|
 \le4^{-n}(34\pi\sqrt3)^nq_{6n}(\Psi).
\end{equation}

For real tests the Gaussian moment estimate is
\[
 \E_P\prod_{a=1}^m|F(f_a)|
 \le c_m\prod_a\|f_a\|_2,\qquad
 c_m=2\,2^{m/2}\Gamma(1+m/2)
 \le2\,2^{m/2}\sqrt{m!};
\]
see \cref{prop:subgaussian} and its proof.
The last inequality is the midpoint log-convexity bound for the
gamma function between \(1\) and \(m+1\).
The \(2n\)-moment estimate on real basis elements shows that
\begin{equation}\label{eq:tensor-random-evaluation}
 Q_\Psi=\sum_{\boldsymbol\kappa}c_{\boldsymbol\kappa}(\Psi)
             \prod_{a=1}^nF(h_{\kappa_a})
\end{equation}
converges absolutely in \(L^2(P)\), with norm at most
\(\sqrt{c_{2n}}\sum|c_{\boldsymbol\kappa}(\Psi)|\).
This proves continuity into \(L^2(P)\).
Expanding each \(f_a\) in the Hermite basis and using the same
product bound identifies the sum with \(\prod_aF(f_a)\) on
decomposable tests. Density of finite Hermite tensors gives
uniqueness of the continuous extension. Finally,
\[
 |S_n(\Psi)|\le c_n\sum_{\boldsymbol\kappa}|c_{\boldsymbol\kappa}(\Psi)|
 \le2\sqrt{n!}\left(\frac{34\pi\sqrt6}{4}\right)^nq_{6n}(\Psi),
\]
which implies \eqref{eq:tensor-growth}.
\end{proof}
